\documentclass[11pt]{article}
\usepackage[T1]{fontenc}
\usepackage{lmodern,amsmath,amssymb,amsthm,mathtools,booktabs,microtype}
\usepackage[margin=1in]{geometry}
\usepackage[hidelinks,pdftitle={Reading the Wake: Period Bounds and Inverse Synthesis for Langton's Ant},pdfauthor={Hao Ke}]{hyperref}
\newtheorem{theorem}{Theorem}
\newtheorem{lemma}[theorem]{Lemma}
\newtheorem{corollary}[theorem]{Corollary}
\newtheorem{proposition}[theorem]{Proposition}
\theoremstyle{remark}\newtheorem{remark}[theorem]{Remark}
\newcommand{\Z}{\mathbb Z}
\newcommand{\N}{\mathbb N}
\newcommand{\R}{\mathrm R}
\newcommand{\Lturn}{\mathrm L}
\newcommand{\ind}{\mathbf 1}
\DeclareMathOperator{\sgn}{sgn}
\title{Reading the Wake:\\Period Bounds and Inverse Synthesis for Langton's Ant}
\author{Hao Ke\\Independent researcher\\\texttt{i@kehao.me}}
\date{15 October 2026}
\begin{document}
\maketitle
\begin{abstract}
A translating periodic trajectory of Langton's ant leaves a wake of permanent
black strands. We show that prescribing this finite set of translation classes,
a non-axis drift and an initial heading determines a computable upper bound on
the period. For diagonal drift, an integer clock gives a signed wake budget,
and a crossing-flow recurrence sharpens it; in particular
$P\le g(3^{W/2}-1)$ for wake growth $g$ and transverse width $W$.
For arbitrary non-axis drift, two reciprocal exponential moments bound $P^2$.
Combining these bounds with the finite-word realization criterion yields a
terminating inverse-synthesis algorithm, with no period supplied as input;
a positive answer constructs a finite black seed. The same budgets exclude
wake-preserving insertions of closed-pose detours. Exact arithmetic checks
and independently verified infinite-trace certificates accompany the proofs.
\end{abstract}

\section{From a wake to a periodic trajectory}
We study the inverse problem of recovering a translating periodic trajectory
from its permanent black wake. The key is to turn that wake into a bound on
the unknown period.

A state is $X=(B,p,h)$, where $B\subset\Z^2$ is finite and $h\in\{N,E,S,W\}$.
At a white cell the ant turns right, at a black cell it turns left, then it flips
the cell and moves one unit. Write $F$ for this map, and observe states before
each update. Put $\chi(x,y)=(-1)^{x+y}$ and
\[
 \sigma(X)=\begin{cases}\chi(p),&h=N,S,\\-\chi(p),&h=E,W.\end{cases}
\]
Both position parity and heading axis change, so $\sigma$ is invariant.

A nonempty turn word $u$ of length $P$, started at $p_0=0$ with specified heading,
defines phase positions $q_i$, headings $h_i$ and a final displacement $d$.
We require heading restoration and $d\ne0$. A realization is a finite black set
whose entire future word is $u^\omega$. The complete black field is \emph{not}
required to translate: a stationary history and a growing wake are permitted.

Partition phase indices by $q_i-q_j\in\Z d$. A class with an odd number of visits
is an \emph{odd class}; write $O\subset\Z^2/\Z d$ for their finite set. For a
realizable word, each odd class contributes one permanent black strand per period.
The input is a drift $d$ with two nonzero coordinates, the initial heading,
and $O$, represented by a finite list of distinct classes. We seek a finite
seed and periodic word realizing exactly these data; \emph{the period is
not supplied}.

The period bound controls all words with the prescribed data, so an exact
word test becomes a terminating search. For the
standard wake with drift $(2,-2)$ the derived cutoff is 856, and the standard
104-step word gives a positive witness.

Jillhewar~\cite{jillhewar} supplies the exact finite-word realization criterion
and constructive seed used here. His width-four analysis gives a special
period--drift bound~\cite{widthfour}. Our contribution is the general
wake-dependent quantitative cutoff: a diagonal crossing recurrence and a
reciprocal-moment bound for every non-axis drift. We use the integer clock
from the companion escape paper~\cite{clock}. After recalling the word test,
we derive both bounds and then the inverse algorithm and detour corollary.

\begin{lemma}[Known finite-word criterion]\label{lem:criterion}
In a translation class choose $b$ and write $q_i=b+a_i d$. Order its phase letters
by increasing $i-Pa_i$. A word is finite-support realizable if and only if every
such ordered word is $\R,\Lturn,\R,\Lturn,\ldots$ starting with $\R$.
\end{lemma}
\begin{proof}
The cell $b+nd$ is read at phase $i$ of epoch $n-a_i\ge0$, at time
$nP+(i-Pa_i)$. For sufficiently large $n$ all class phases are present and the
cell is initially white, since the seed is finite. Its consecutive letters must
alternate, proving necessity. For sufficiency, if
$\min a_i\le n<\max a_i$, initially blacken $b+nd$ precisely when the number of
removed phases, those with $a_i>n$, is odd. All other cells start white.
The available phase word is a suffix of the stabilized alternating word and
starts with the prescribed colour. Induction on physical time proves the whole
infinite trace. Only finitely many boundary cells are blackened.
\end{proof}
In particular $\#\R-\#\Lturn=|O|=:g$. Heading restoration implies $g\equiv0$
modulo four, $P$ even and $d_x+d_y$ even. The quotient is by the \emph{full} drift:
for $d=(m,-m)$ a class needs both $x+y$ and $x\bmod m$.

\section{An integer clock for diagonal budgets}
After rotating a diagonal drift parallel to $(1,-1)$, the transverse
coordinate $x+y$ and heading return after each period. A clock using these
quantities has a cancelling
head correction: its period gain comes from the wake. We recall the unit
increment identity of~\cite{clock} before computing that gain.

Fix an integer $c$, set $s=x+y-c$, and put
\[
 f(s)=3^{|s|}-1,\qquad
 \epsilon(h)=\begin{cases}+1,&h=N,E,\\-1,&h=S,W,\end{cases}\qquad
 G(s,e)=\begin{cases}\tfrac12 f(s),&es\ge0,\\-\tfrac32 f(s),&es<0.\end{cases}
\]
Use the integer clock of~\cite{clock} on the complete finite state
\begin{equation}\label{eq:integer}
 \Theta_c(X)=-\sigma(X)\sum_{q\in B}\chi(q)f(q_x+q_y-c)
             +G(p_x+p_y-c,\epsilon(h)).
\end{equation}
Since $f(s)$ is even, $\Theta_c$ is an exactly computable integer.

\begin{lemma}[Integer clock, Ke~\cite{clock}]\label{thm:clock}
For every finite-support state and integer centre $c$, with $e,e'$ the before
and after values of $\epsilon$ and $s=p_x+p_y-c$,
\begin{equation}\label{eq:gain}
 \Theta_c(FX)-\Theta_c(X)
   =1+4f(s)\ind_{\{e'\ne e,\ es<0\}}\ \ge1.
\end{equation}
\end{lemma}
\begin{proof}
Write $\delta=+1$ on $\R$ and $-1$ on $\Lturn$. Directly checking the two colours
at four headings gives $\delta\sigma\chi(p)=ee'$, and the outgoing step gives
$s'=s+e'$. Thus the black-field increment is $-ee'f(s)$. Let $G_\pm(s)=G(s,\pm1)$.
The elementary recurrences are
\begin{align*}
 G_+(s+1)-G_+(s)&=f(s)+1,&G_-(s-1)-G_-(s)&=f(s)+1,\\
 G_-(s)-G_+(s)&=-2\sgn(s)f(s).
\end{align*}
If $e'=e$ the total gain is one. For $+1\to-1$ it is
$1+2f(s)+G_-(s)-G_+(s)$, equal to one for $s\ge0$ and to $1+4f(s)$ for $s<0$.
The other reversal gives the reflected cases. This includes $s=0$.
\end{proof}
\section{Diagonal wake budgets and crossing refinement}
Rotate a diagonal drift to $d=(m,-m)$, $m>0$, and use the transformed heading.
Let $t=x+y$, let $o_t$ count odd classes at level $t$, and define
\begin{equation}\label{eq:A}
 A_t=-\sigma(-1)^t o_t,\qquad
 K_c(O)=\sum_t A_t(3^{|t-c|}-1).
\end{equation}
\begin{theorem}[Wake-determined diagonal period]\label{thm:diagonal}
Every realization of this wake obeys $P\le K_c(O)$ for every integer $c$.
\end{theorem}
\begin{proof}
Over one period the head term of \eqref{eq:integer} cancels, since both $t$ and
the heading return to their initial values. The signed cell toggles give
$-\sigma\sum_i\chi(q_i)f(t_i-c)\delta_i$.
The weight is constant in each full $d$-class. Lemma~\ref{lem:criterion} makes
its net $\delta$ count one on odd classes and zero on even classes. The exact
period increment is therefore $K_c(O)$. Summing \eqref{eq:gain} proves the bound.
The unchanged part of the full field cancels algebraically; it is not deleted.
\end{proof}

Here is an independent geometric derivation, which also supplies a refinement.
The finite projected period is a closed walk with increments $\pm1$ in $t$.
Let $f_t$ count upward crossings of the cut $(t,t+1)$; it also counts downward
crossings. Let $m_t$ count visits to $t$, and $u_t$ count passages entering and
leaving in the increasing direction. Flow conservation equates this with the
number of decreasing-through passages. The fixed arrival axis and the two
turn choices give
\begin{equation}\label{eq:flow}
 m_t=f_{t-1}+f_t=A_t+4u_t,\qquad
 0\le u_t\le\min(f_{t-1},f_t).
\end{equation}
For example, at vertical-arrival levels, right turns are through passages and
left turns reverse the projection, giving $\#\R-\#\Lturn=4u_t-m_t$; horizontal
arrival reverses that sign. Substituting $u_t\le f_{t-1}$ into
$f_t=A_t+4u_t-f_{t-1}$ gives the first inequality below; $u_t\le f_t$
gives the second:
\begin{equation}\label{eq:recurrence}
 f_t\le3f_{t-1}+A_t,\qquad f_{t-1}\le3f_t+A_t.
\end{equation}
Let $\ell,h$ be the extreme wake levels. Outside them $A_t=0$; starting past the
finite support of the projected walk and applying \eqref{eq:recurrence} inward
forces $f_{\ell-1}=f_h=0$. Thus the periodic trace lies between the extreme wake levels. An empty wake forces every
$f_t=0$, impossible for a nonempty period.

Starting from these two zero boundary values, set
\[
 L_{\ell-1}=0,\quad L_t=3L_{t-1}+A_t,\qquad
 R_h=0,\quad R_t=3R_{t+1}+A_{t+1},\quad C_t=\min(L_t,R_t).
\]
Then $0\le f_t\le C_t$ for $\ell\le t<h$. If a $C_t$ is negative, the wake is
impossible; otherwise
\begin{equation}\label{eq:cuts}
 P=2\sum_t f_t\le2\sum_t C_t.
\end{equation}
Moreover $d_x=\sigma\sum_t(-1)^t f_t$: on a cut with $\sigma(-1)^t=1$, the
upward horizontal moves contribute $+f_t$ to $x$; on the other cuts the downward
horizontal moves contribute $-f_t$. If $C_+$ and $C_-$ sum the corresponding
positive and negative cut bounds, then also
\begin{equation}\label{eq:direction}
 P\le4C_+-2d_x,\qquad P\le4C_-+2d_x.
\end{equation}

\begin{corollary}[Width and growth bound]
Writing $W=h-\ell$, a diagonal realization has even $W$ and
\begin{equation}\label{eq:coarse}
 P\le g(3^{W/2}-1).
\end{equation}
\end{corollary}
\begin{proof}
At a transverse extreme every passage reverses the projection. The fixed arrival
axis makes all letters in an extremal class equal. Lemma~\ref{lem:criterion}
therefore forces exactly one $\R$ per such class. Both extremes have the same
arrival axis, so $W$ is even; this is the usual extremal rigidity argument.
Take $c=(\ell+h)/2$ in Theorem~\ref{thm:diagonal}, discard negative coefficients
and use $\sum_t\max(A_t,0)\le g$ to obtain \eqref{eq:coarse}.
\end{proof}

For a phase of standard P104 with normalized drift $(2,-2)$ and $\sigma=1$,
the coefficients $A_t$, for $t=-3,\ldots,7$, and the cut bounds are
\begin{align*}
 A&=(2,-1,1,-1,1,-1,1,0,1,-1,2),\\
 C&=(2,5,16,47,142,145,48,16,5,2).
\end{align*}
Their cut sum is $428$, so \eqref{eq:cuts} gives $P\le856$.
Equation~\eqref{eq:A} with $c=2$ gives the same value. The standard word
certifies existence within this finite search domain; classifying every
realization in the domain remains an explicit computational problem.

\section{General non-axis drift by reciprocal moments}
For a general drift, use a transverse linear coordinate constant along it.
An exponential weight and its reciprocal give complementary period budgets;
Cauchy--Schwarz then bounds the number of steps independently of their
individual weights.

Rotate the drift to $(a_0,-b_0)$, $a_0,b_0>0$. Put
$k=\gcd(a_0,b_0)$, $\eta=(b_0/k,a_0/k)$, and choose rational $r>1$ such that,
with $a=r^{\eta_x}$ and $b=r^{\eta_y}$,
\[
 \kappa=1-\frac2{a+1}-\frac2{b+1}>0.
\]
The choice $r=4$ always works. Define $w(q)=r^{\eta\cdot q}$, and
\[
 (\gamma_N,\gamma_E,\gamma_S,\gamma_W)
 =\left(\frac2{b+1},\frac2{a+1},-\frac{2b}{b+1},-\frac{2a}{a+1}\right).
\]
The one-sided potentials are
\[
 \Psi_+(X)=-\sigma\sum_{q\in B}\chi(q)w(q)+\gamma_h w(p),\qquad
 \Psi_-(X)=-\sigma\sum_{q\in B}\frac{\chi(q)}{w(q)}+\frac{\gamma_{h+2}}{w(p)}.
\]
Substitution of the current-cell flip gives $\Delta\Psi_+=\kappa w(p)$ in six
heading/colour cases and $(\kappa+4)w(p)$ in $S/\Lturn$ and $W/\R$.
The half-turn formula gives $\Delta\Psi_-\ge\kappa/w(p)$. These local identities
are the one-sided exponential clock construction, included to make the new
period argument self-contained.

Define the two finite wake moments, using any representatives of $O$,
\[
 A_+=-\sigma\sum_{q\in O}\chi(q)r^{\eta\cdot q},\qquad
 A_-=-\sigma\sum_{q\in O}\chi(q)r^{-\eta\cdot q}.
\]
They are independent of the representatives since $\eta\cdot d=0$ and
$d_x+d_y$ is even. The head terms cancel over a period, and
Lemma~\ref{lem:criterion} gives exact increments $A_+$ and $A_-$. Thus both are
positive and $\kappa\sum_i w(q_i)\le A_+$,
$\kappa\sum_i w(q_i)^{-1}\le A_-$. Cauchy--Schwarz yields:
\begin{theorem}[Reciprocal wake bound]\label{thm:moment}
Every realization with the specified non-axis drift and wake satisfies
\begin{equation}\label{eq:moment}
 P^2\le \frac{A_+A_-}{\kappa^2}.
\end{equation}
\end{theorem}
All quantities are rational. The integer square root of the floor of the right
side is a computable cutoff. Changing the transverse centre scales $A_+$ and
$A_-$ reciprocally, leaving the product unchanged. For axis drift the normal has
a zero coordinate, so this construction cannot give positive $\kappa$.

\section{All-period synthesis and detour obstruction}
\begin{theorem}[Period-unknown inverse synthesis]\label{thm:inverse}
Given a non-axis integer drift, an initial heading and a finite odd-class wake,
existence of a finite-support periodic realization is decidable. A positive
answer can produce its word and an explicit finite seed.
\end{theorem}
\begin{proof}
First reject odd $d_x+d_y$, growth $|O|$ not divisible by four, or a
nonpositive moment $A_+$ or $A_-$. Otherwise compute
Theorem~\ref{thm:moment}'s cutoff, additionally using
Theorem~\ref{thm:diagonal} and \eqref{eq:cuts}--\eqref{eq:direction} in the
diagonal case. A negative crossing bound also rejects the input.
Enumerate every nonempty word up to the least derived bound.
Check its heading and drift, apply Lemma~\ref{lem:criterion}, and require the
exact odd-class set. A passing word constructs a finite seed; complete finite
exhaustion gives a negative answer. Completeness follows from the period bound.
\end{proof}
The implementation checks constructed seeds by a second, all-future formula.
At phase $i$, epoch $n$, it evaluates
\[
 \ind_B(q_i+nd)\ \mathbin{\mathrm{xor}}\
 \left|\{(j,s):q_j-q_i=sd,\ s>0\text{ or }(s=0,j<i),\ s\le n\}\right|\bmod2.
\]
Only finitely many initial black hits and previous-write thresholds can change
this value. Checking all breakpoints also checks the last infinite interval.
An optional execution cap reports \texttt{INCONCLUSIVE}, not a negative answer.
The exhaustive search has up to $2^{B+1}-2$ candidate words for cutoff $B$.

\begin{corollary}[No wake-preserving closed-pose insertion]
Suppose realizable words $u,v$ have the same non-axis drift, and $v$ is obtained
by inserting nonempty detours into $u$, each returning to the same position and
heading while leaving the original steps unchanged. Then
\[
 A_+(O_v)>A_+(O_u),\qquad A_-(O_v)>A_-(O_u).
\]
Thus their signed transverse wake cannot be identical. For diagonal drift,
$K_c(O_v)-K_c(O_u)\ge |v|-|u|>0$ for every integer $c$.
\end{corollary}
\begin{proof}
Every retained original step has the same position, heading and letter; hence
it has the same local clock increment. Every inserted step has strictly positive
increment in both one-sided clocks, and at least one in the integer clock.
Sum over the words and use the exact wake-increment equalities.
\end{proof}
Thus any such realizable detour requires an additional wake budget.

\section{Exact verification and examples}
The exact checker tests integer and rational increments, projected flow and
full-period moment identities. All 104 standard phases, four rotations and
powers one through three satisfy the $856$-scaled cutoff, giving 1,248 positive
controls. Separate polynomial audits verify the local identities for arbitrary
parameters and all integer-clock sign branches. The all-future breakpoint
test verifies complete seeds, including remote initial black cells.

For a known excluded-width control with drift $(1,-1)$ and heading $N$, take wake
representatives
\[ (0,-3),\quad (0,-1),\quad (0,0),\quad (0,1). \]
The derived cutoff is 14.
The pruned search exhausts 2,226 nodes and returns a negative answer; a separate
enumeration of all 32,766 nonempty words through length 14 agrees.
The standard wake supplies a positive witness; the full search through
its cutoff 856 remains open.

A supplementary experiment inserts all 14,386 internally alternating
closed-pose fragments of lengths at most 28 at all 104 standard boundaries.
The 1,496,144 insertion pairs produce no additional realizable word.
An independent Python search reproduces the complete smaller domain through
length 20.

\section{Conclusion}
A finite wake carries a quantitative budget for the period that produces it.
The diagonal integer budget and crossing recurrence expose this relation
combinatorially, while reciprocal exponential moments extend it to every
non-axis drift. Together with exact finite-word realization, these bounds
turn the prescribed-wake problem into a finite synthesis procedure.
Extracting possible wakes from arbitrary initial states, treating axis drift
and classifying the standard wake are further questions.

\paragraph{Reproducibility.}
The manuscript, proofs, exact checkers and retained reports are in the
\href{https://github.com/kehao95/langtons-ant-core/tree/main/wake_synthesis}{wake synthesis project}.
From the repository root, \texttt{python3 wake\_synthesis/reproduce.py} runs
the main checks using Python's standard library. Adding \texttt{--search}
runs the finite detour experiment and requires a C++17 compiler.
Build the paper with \texttt{make -C wake\_synthesis/paper}.
The release manifest records source hashes and verification results.

\paragraph{AI assistance.}
Generative AI assisted with developing and checking arguments, implementing
verification programs, and revising the exposition.

\begin{thebibliography}{9}
\bibitem{jillhewar}
A. Jillhewar, \emph{Finite-Support Periodic Highways of Langton's Ant:
Necessary Conditions, Transverse Exclusions, and Exact Search}, preprint, 2026.
\href{https://doi.org/10.21203/rs.3.rs-10516269/v1}{doi:10.21203/rs.3.rs-10516269/v1}.

\bibitem{widthfour}
A. Jillhewar, \emph{Width-four research log}, 21 July 2026, accompanying source
note, \href{https://github.com/Atharva12456/langtons-ant-highway/blob/76a7e234e577d992d086c204c3dea00ae72cf10a/research-notes/width4_research_2026-07-21.md}{revision \texttt{76a7e234e577}}.

\bibitem{clock}
H. Ke, \emph{An Ant on the Clock: Exponential Escape Bounds for Langton's Ant},
Zenodo, version 0.4.0, 2026.
\href{https://doi.org/10.5281/zenodo.23402379}{doi:10.5281/zenodo.23402379}.

\end{thebibliography}
\end{document}
